\section*{Chapter \ref{ch:iv:machine-learning} --- Machine Learning}

\begin{solution}{iq:iv:machine-learning:1}
More data: variance falls, bias unchanged. More features or flexibility: bias falls, variance rises. More regularisation:
bias rises, variance falls. With noisy financial targets the variance term usually dominates, so the gains come from data
and regularisation, not flexibility.
\iqlookfor{the three directions, and which term dominates in finance.}
\end{solution}

\begin{solution}{iq:iv:machine-learning:2}
The constraint region of the lasso ($\|\beta\|_1 \le t$) has corners on the axes, and the elliptical contours of the squared
loss usually first touch it at a corner, where some coordinates are zero. The ridge region is a ball with no corners, so the
contact point generically has all coordinates non-zero. Equivalently, the lasso's subgradient at zero is an interval, so zero
is optimal for small effects.
\iqlookfor{the geometric or subgradient explanation.}
\end{solution}

\begin{solution}{iq:iv:machine-learning:3}
It depends on the payoffs and on the evidence. If gains and losses are symmetric and unrelated to the prediction, 55\% is a
large edge for daily data; if the model is right on small moves and wrong on large ones, it can lose money. Ask for the
expected P\&L per prediction, a proper score of the probabilities (log loss or Brier), the sample size (55\% over 250 days has a
standard error of about 3 points) and how many models were tried.
\iqlookfor{payoff asymmetry, uncertainty of the estimate and the right metric.}
\end{solution}

\begin{solution}{iq:iv:machine-learning:4}
The penalty treats all coefficients alike, so features on larger scales are penalised less per unit of effect; standardising
puts them on equal footing. Use the mean and standard deviation of the training fold only (fitted inside each validation
split), otherwise the scaling leaks information from the test period.
\iqlookfor{scale dependence of the penalty and leak-free preprocessing.}
\end{solution}

\begin{solution}{iq:iv:machine-learning:5}
With nearly collinear features the lasso's solution is unstable: small changes in the sample flip which feature carries the
effect (\cref{ex:iv:machine-learning:lasso}: one coefficient zeroed in 63\% of samples, either feature equally often). Report
the pair as one effect, use ridge or an elastic net to share the weight, or combine the features; never interpret the lasso's
choice between them.
\iqlookfor{instability of sparse selection under collinearity and a robust alternative.}
\end{solution}

\begin{solution}{iq:iv:machine-learning:6}
Leakage through overlapping labels and shuffled folds: consecutive days have almost the same features and almost the same
20-day target, so a test day's neighbours in the training set give away its label. Test it on data with no signal: on the
chapter's noise panels the same pipeline scores 0.26 with shuffled folds and $-0.61$ with purged walk-forward evaluation
(\cref{fig:iv:machine-learning:leak}). Re-evaluate with walk-forward splits, purge training labels that overlap the test
period, and add an embargo.
\iqlookfor{the overlap mechanism and a noise-panel test of the pipeline.}
\end{solution}

\begin{solution}{iq:iv:machine-learning:7}
It fitted noise: the gap between training and test $R^2$ is the variance of an over-flexible model (the chapter's
demonstration uses pure noise and gets these numbers). Use fewer, shallower trees with a lower learning rate, subsampling, and
early stopping on a time-ordered validation set; compare with a linear baseline; and check that the test set is forward in
time.
\iqlookfor{recognising overfitting and the standard boosting controls.}
\end{solution}

\begin{solution}{iq:iv:machine-learning:8}
Impurity importance counts how much each split reduces training impurity, and a continuous noisy feature offers many split
points, so the forest finds spurious splits on it; the importance is computed in sample. Use permutation importance on
held-out, time-ordered data (the drop in the validation metric when the feature is shuffled), with repeats to measure its noise.
\iqlookfor{the bias of impurity importance and a held-out alternative.}
\end{solution}

\begin{solution}{iq:iv:machine-learning:9}
Ten years are about 2\,500 dependent observations of a target whose predictable part is a fraction of a percent of its
variance; a network with many parameters fits the noise, and its training is unstable across seeds. Non-stationarity makes old
data a poor guide. A regularised linear model has little variance and captures most of what is there. Networks help where data
are plentiful and structured (order books, text), not on daily returns.
\iqlookfor{sample size against signal strength, and non-stationarity.}
\end{solution}

\begin{solution}{iq:iv:machine-learning:10}
Decision: where and whether to post, trading off fill probability against adverse selection. Target: filled within one second
of placement (with partial fills defined), observed from our own order records. Features known at placement: queue position
and size ahead, spread, depth on both sides, recent trade rate and imbalance, time of day, volatility, the order's price
relative to the touch. Baseline: a logistic regression; then boosting. Validate forward in time on our own orders only, beware
of selection (we only observe outcomes for orders we sent), and calibrate the probabilities (proper scoring rule). Monitor
calibration drift and the feedback loop (the model changes where we post, which changes the data).
\iqlookfor{the full chain from decision to monitoring, including selection bias and feedback.}
\end{solution}

\begin{solution}{iq:iv:machine-learning:11}
Define toxicity as the mark-out: the move of the mid price against us after the counterparty's trade, over horizons such as one
second and one minute. Target: the mark-out of each fill (regression) or its sign beyond a threshold. Features by counterparty and
order: past mark-outs of the counterparty, order size, timing relative to news or other venues' moves, the counterparty's order-to-
trade ratio. Validate by counterparty and forward in time; the metric is the P\&L improvement from widening on the model's flags,
not accuracy. Watch fairness and regulatory constraints on discriminating between clients, and the feedback when toxic clients
adapt.
\iqlookfor{mark-out as the target, counterparty-level validation and a P\&L metric.}
\end{solution}

\begin{solution}{iq:iv:machine-learning:12}
In order: (1) leakage through overlapping labels or shuffled folds: re-run with purged walk-forward evaluation; (2) look-ahead in
features (a close used before it was known, restated data): audit each feature's timestamp; (3) selection over many models or
parameters: count the trials and deflate; (4) regime change: compare feature distributions live and in sample; (5) execution:
the live prediction arrives later than assumed. A 71\% accuracy on daily direction is itself implausible, which is the first
clue.
\iqlookfor{a ranked list of causes, each with a concrete check.}
\end{solution}

\begin{solution}{iq:iv:machine-learning:13}
Trade when $p \times 1 > (1 - p) \times 2$, that is $p > \tfrac23$. Accuracy weights all errors alike and uses a threshold of one
half, so a model with high accuracy can lose money; evaluate the expected P\&L of the decision rule, and score the probabilities
with a proper rule (log loss) so that they can be thresholded correctly.
\iqlookfor{the payoff-based threshold and the case for proper scoring.}
\end{solution}
